% GENERATED FILE. Edit canonical lessons, then run npm run generate:books.
% canonical-source-hash: fnv1a64:5eeb6e40797c1fa9
% canonical-lessons: UR-C137-jamni, UR-C137-sunahra, UR-C137-band, UR-C137-khula, UR-C137-jagna

\chapter{Purple, Golden, Closed, Open, To wake up}
\label{ch:ur-purple-golden-closed-open-to-wake-up}
\hlchaptermodality{137}

\begin{chapteropening}
\textbf{By the end of this chapter:} \emph{I can say purple, golden, closed, open and to wake up.}

Complete the last lesson of chapter 137: I can say purple, golden, closed, open and to wake up.
\end{chapteropening}

\section[\texorpdfstring{\ur{جامنی} (jāmnī)}{جامنی (jāmnī)}]{\ur{جامنی} (jāmnī) — purple}

\label{lesson:UR-C137-jamni}



\begin{quote}
Before the new one: say the Urdu for ripe, then the Urdu for pink.
\end{quote}

\subsection*{You'll want to know: \ur{جامنی}}
\textbf{\ur{جامنی}} — \emph{jāmnī} — \textquotedblleft{}purple\textquotedblright{}.

\begin{cousinweb}
\textbf{\ur{جامنی}} (\emph{jāmnī}) is \textquotedblleft{}jāmun-coloured\textquotedblright{}, the purple of the jambolan fruit, with the ending \emph{-ī}. CDIAL lists Hindi \emph{jāmun} and \emph{jāman}, the tree and its fruit, under \emph{*jambūna}, a form beside Sanskrit \emph{jambula} (CDIAL 5136).
\end{cousinweb}

\begin{grammarlens}[title={What you've built}]
One more word for this chapter.
\end{grammarlens}

\subsection*{Your turn}
\begin{itemize}
\raggedright
  \item \emph{Say it:} \emph{jāmnī}
  \item \emph{Say it:} \emph{jāmnī}, once more
  \item \emph{Say it:} say \emph{jāmnī}
  \item \emph{Recall:} say the Urdu for ripe, then the Urdu for pink, then say \emph{jāmnī} again
\end{itemize}

\begin{tcolorbox}[breakable,colback=teal!4,colframe=teal!35!black,title={Before you move on}]
What does \ur{جامنی} mean? (\textquotedblleft{}Purple\textquotedblright{}.) Say it once more.
\end{tcolorbox}
\section[\texorpdfstring{\ur{سنہرا} (sunahrā)}{سنہرا (sunahrā)}]{\ur{سنہرا} (sunahrā) — golden}

\label{lesson:UR-C137-sunahra}



\begin{quote}
Before the new one: say the Urdu for pink, then the Urdu for purple.
\end{quote}

\subsection*{You'll want to know: \ur{سنہرا}}
\textbf{\ur{سنہرا}} — \emph{sunahrā} — \textquotedblleft{}golden\textquotedblright{}.

\begin{cousinweb}
\textbf{\ur{سنہرا}} (\emph{sunahrā}) is a tadbhava, a word inherited and reshaped by sound change. It comes from \emph{*suvarṇadhara}, \textquotedblleft{}containing gold\textquotedblright{}, a form CDIAL reconstructs from its descendants (CDIAL 13522). It is built from \emph{suvarṇa}, \textquotedblleft{}gold\textquotedblright{}, and \emph{dhara}, \textquotedblleft{}holding\textquotedblright{}.
\end{cousinweb}

\begin{grammarlens}[title={What you've built}]
One more word for this chapter.
\end{grammarlens}

\subsection*{Your turn}
\begin{itemize}
\raggedright
  \item \emph{Say it:} \emph{sunahrā}
  \item \emph{Say it:} \emph{sunahrā}, once more
  \item \emph{Say it:} say \emph{sunahrā}
  \item \emph{Recall:} say the Urdu for pink, then the Urdu for purple, then say \emph{sunahrā} again
\end{itemize}

\begin{tcolorbox}[breakable,colback=teal!4,colframe=teal!35!black,title={Before you move on}]
What does \ur{سنہرا} mean? (\textquotedblleft{}Golden\textquotedblright{}.) Say it once more.
\end{tcolorbox}
\section[\texorpdfstring{\ur{بند} (band)}{بند (band)}]{\ur{بند} (band) — closed}

\label{lesson:UR-C137-band}



\begin{quote}
Before the new one: say the Urdu for purple, then the Urdu for golden.
\end{quote}

\subsection*{You'll want to know: \ur{بند}}
\textbf{\ur{بند}} — \emph{band} — \textquotedblleft{}closed\textquotedblright{}.

\begin{cousinweb}
\textbf{\ur{بند}} (\emph{band}) is a loan, usually traced to Persian \emph{band}, \textquotedblleft{}shut, bound\textquotedblright{}.
\end{cousinweb}

\begin{grammarlens}[title={What you've built}]
One more word for this chapter.
\end{grammarlens}

\subsection*{Your turn}
\begin{itemize}
\raggedright
  \item \emph{Say it:} \emph{band}
  \item \emph{Say it:} \emph{band}, once more
  \item \emph{Say it:} say \emph{band}
  \item \emph{Recall:} say the Urdu for purple, then the Urdu for golden, then say \emph{band} again
\end{itemize}

\begin{tcolorbox}[breakable,colback=teal!4,colframe=teal!35!black,title={Before you move on}]
What does \ur{بند} mean? (\textquotedblleft{}Closed\textquotedblright{}.) Say it once more.
\end{tcolorbox}
\section[\texorpdfstring{\ur{کھلا} (khulā)}{کھلا (khulā)}]{\ur{کھلا} (khulā) — open}

\label{lesson:UR-C137-khula}



\begin{quote}
Before the new one: say the Urdu for golden, then the Urdu for closed.
\end{quote}

\subsection*{You'll want to know: \ur{کھلا}}
\textbf{\ur{کھلا}} — \emph{khulā} — \textquotedblleft{}open\textquotedblright{}.

\begin{cousinweb}
\textbf{\ur{کھلا}} (\emph{khulā}) is a tadbhava, a word inherited and reshaped by sound change. It comes from \emph{*khull}, \textquotedblleft{}to be open\textquotedblright{}, a form CDIAL reconstructs from its descendants (CDIAL 3945).
\end{cousinweb}

\begin{grammarlens}[title={What you've built}]
One more word for this chapter.
\end{grammarlens}

\subsection*{Your turn}
\begin{itemize}
\raggedright
  \item \emph{Say it:} \emph{khulā}
  \item \emph{Say it:} \emph{khulā}, once more
  \item \emph{Say it:} say \emph{khulā}
  \item \emph{Recall:} say the Urdu for golden, then the Urdu for closed, then say \emph{khulā} again
\end{itemize}

\begin{tcolorbox}[breakable,colback=teal!4,colframe=teal!35!black,title={Before you move on}]
What does \ur{کھلا} mean? (\textquotedblleft{}Open\textquotedblright{}.) Say it once more.
\end{tcolorbox}
\section[\texorpdfstring{\ur{جاگنا} (jāgnā)}{جاگنا (jāgnā)}]{\ur{جاگنا} (jāgnā) — to wake up}

\label{lesson:UR-C137-jagna}



\begin{quote}
Before the new one: say the Urdu for closed, then the Urdu for open.
\end{quote}

\subsection*{You'll want to know: \ur{جاگنا}}
\textbf{\ur{جاگنا}} — \emph{jāgnā} — \textquotedblleft{}to wake up\textquotedblright{}.

\begin{cousinweb}
\textbf{\ur{جاگنا}} (\emph{jāgnā}) is a tadbhava, a word inherited from Sanskrit and reshaped by sound change. It comes from \emph{jāgrati}, \textquotedblleft{}are awake\textquotedblright{} (CDIAL 5175).
\end{cousinweb}

\begin{grammarlens}[title={What you've built}]
One more word for this chapter.
\end{grammarlens}

\subsection*{Your turn}
\begin{itemize}
\raggedright
  \item \emph{Say it:} \emph{jāgnā}
  \item \emph{Say it:} \emph{jāgnā}, once more
  \item \emph{Say it:} say \emph{jāgnā}
  \item \emph{Recall:} say the Urdu for closed, then the Urdu for open, then say \emph{jāgnā} again
\end{itemize}

\begin{tcolorbox}[breakable,colback=teal!4,colframe=teal!35!black,title={Before you move on}]
What does \ur{جاگنا} mean? (\textquotedblleft{}To wake up\textquotedblright{}.) Say it once more.
\end{tcolorbox}
